\documentclass[10pt,a4paper]{amsart}
\usepackage[margin=19mm]{geometry}
\usepackage{amsmath,amssymb,mathtools}
\usepackage[scr=rsfs,cal=euler]{mathalfa}
\usepackage{xfrac}
\usepackage[unicode,hidelinks,bookmarks=false]{hyperref}
\newcommand{\ie}{i.e.\ }
\newcommand{\cf}{cf.\ }
\newcommand{\resp}{respectively\ }
\newcommand{\Subsection}{\subsection}
\providecommand{\nobreakdash}{\nobreak\hbox{-}}
\setlength{\emergencystretch}{2em}
\multlinegap0pt
\usepackage{enumerate}
\usepackage{amssymb}
\usepackage{mathtools}
\usepackage{dsfont}
\usepackage[normalem]{ulem}
\newtheorem{theorem}{Theorem}
\newcommand{\EE}{\mathbb{E}}
\newcommand{\NN}{\mathbb{N}}
\newcommand{\RR}{\mathbb{R}}
\title{Collisions of random walks in unimodular random graphs: applications to the random geometric graph and long-range percolation}
\author{Jhon Astoquillca, Carlos Martinez-Arevalo}
\date{Source: arXiv:2607.04002v1 (CC BY 4.0)}
\begin{document}
\maketitle

\noindent\emph{Source and licence.} Collisions of random walks in unimodular random graphs: applications to the random geometric graph and long-range percolation. Version arXiv:2607.04002v1 (CC BY 4.0), distributed by arXiv under the Creative Commons Attribution 4.0 International licence, http://creativecommons.org/licenses/by/4.0/, as recorded in the arXiv licence metadata for this exact version and shown on the abstract page https://arxiv.org/abs/2607.04002v1. Full licence notice: \texttt{licenses/arxiv-2607.04002-CC-BY-4.0.txt}.\par\smallskip

\noindent\textit{Editorial scope.} Counted unit: the article's theorem thm\_unimodularity\_FCP (source0.tex lines 107--109). The article's complete original proof of it is delivered with it (source0.tex lines 587--610, one \texttt{proof} environment of 24 lines, which the article places away from the statement). No mathematics is written, repaired or re-derived: the statement and the proof are the article's own lines, in the article's own notation, and no retained line is substituted or re-flowed.  Retained uncounted as the source-local context the proof needs: lines 531--533; lines 583--585.  Nothing was omitted from any retained span, so the package carries no omission record.  No retained line contains a citation, so the wrapper carries no bibliography and no bibliography entry.  No retained line contains a figure environment, an image-inclusion command or drawing code, so no image is delivered and the package declares no asset dependency.  Editorial formatting only, in addition: an AMS \texttt{amsart} preamble that restates the article's own declarations and macros used by the retained lines, namely the environments \texttt{theorem} on separate counters, together with the standard packages those lines need.  The retained environments print in this document as Theorem 1 in order of appearance, whereas the article prints the same items with the numbering of its own full document.  The complete native source of the article as distributed in the arXiv e-print for this exact version is bundled unchanged as \texttt{sources/204438/source0.tex}.
\begin{equation}\label{Mass_Transport_Principle}
    \mathbf E \Big[\sum_{v \in V} f(G,\rho,v) \Big] = \mathbf E \Big[\sum_{u \in V} f(G,u,\rho) \Big].
\end{equation}

\begin{equation}\label{eq_standar_cont_disc}
\int^\infty_0 p^\mathrm{cont}_t(x,x) \; \mathrm{d}t = \sum_{n \ge 0}p_n(x,x), \quad \text{for all } x \in V.
\end{equation}

\begin{theorem}\label{thm_unimodularity_FCP}
Let~$(G,\rho)$ be a unimodular random rooted graph with~$\mathbf E[ \deg(\rho) \cdot \sum_{n \ge 0} p_n(\rho,\rho) ] <  \infty$. Then~$G$ has the discrete and continuous finite collision property almost surely. 
\end{theorem}

\begin{proof}[Proof of Theorem~\ref{thm_unimodularity_FCP}] We assume that~$(G,\rho)$ is unimodular and that~$\mathbf{E} [\deg(\rho) \cdot \mathsf G(\rho)] < \infty$. We denote by~$\mathbf P^G$ the joint law of two independent discrete-time random walks and two independent continuous-time random walks on~$G$, and by~$\mathbf E^G$ its expectation operator. We define
$$ Z_\mathrm{disc}(x) := |\{n \in \NN: X_n = Y_n\}| \quad \text{and} \quad Z_\mathrm{cont}(x) := |\{ t \ge 0: X^\mathrm{cont}_{t^-} \neq Y^\mathrm{cont}_{t^-}, \; X^\mathrm{cont}_t = Y^\mathrm{cont}_t \}| $$
with the random walks started both from~$x \in V(G)$. 

We start by proving that~$G$ has the discrete finite collision property almost surely. Using the reversibility of the random walk, we have 
\begin{equation}\label{eq_pre_apply_MTP}
    \deg(x) \cdot \mathbf E^G [Z_\mathrm{disc}(x)] = \sum^\infty_{n = 0} \sum_{v \in V(G)} p_n(x,v) \cdot p_n(v,x) \cdot \deg(v),
\end{equation}
Consider the mass transport function
$$ f(G,u,v) = p_n(u,v) \cdot p_n(v,u) \cdot \deg(v) $$
Each vertex~$u$ sends a total mass of~$\sum_{v \in V(G)}p_n(u,v) p_n(v,u) \deg(v)$, while, each vertex~$v$ receives a total mass of
$$ \sum_{u \in V(G)} f(G,u,v) = \deg(v) \cdot p_{2n}(v,v). $$
We apply the Mass-Transport Principle~\eqref{Mass_Transport_Principle} in~\eqref{eq_pre_apply_MTP} to obtain that 
$$ \mathbf{E} \big[ \deg(\rho) \cdot \mathbf E^G[Z_\mathrm{disc}(\rho) ] \big] = \mathbf{E} \big[ \deg(\rho) \cdot \sum^\infty_{n = 0} p_{2n}(\rho,\rho) \big] \leq \mathbf E[ \deg(\rho) \cdot \mathsf G(\rho,\rho) ]. $$
Hence, the hypothesis implies that~$Z_\mathrm{disc}(\rho)<\infty$~$\mathbf P$-a.s., which, combined with the connectedness of~$G$, yields that~$G$ has the discrete finite collision property almost surely.

We now prove the continuous analogue. To this end, we use Wald's equation and then Fubini's theorem to see that
$$ 2 \cdot \mathbf E^G [ Z_\mathrm{cont}(x)] = \EE^G[ \mathrm{Leb}(t \ge 0: X^\mathrm{cont}_t = Y^\mathrm{cont}_t) ] = \int^\infty_0 \sum_{v \in V(G)} p^\mathrm{cont}_t(x,v) \cdot p^\mathrm{cont}_t(v,x) \cdot \frac{\deg(v)}{\deg(x)}, $$
where Leb stands for the Lebesgue measure on~$\RR$. Applying the Mass–Transport Principle in the same way as in the discrete case, we obtain that
\begin{equation*}
    \mathbf E \big[ 2 \deg(\rho) \cdot \mathbf E^G[Z_\mathrm{cont}(\rho)] \big] = \mathbf E \big[ \deg(\rho) \cdot \int^\infty_0 p^\mathrm{cont}_{2t}(\rho,\rho) \; \mathrm{d}t \big].
\end{equation*}
Finally, changing variables~$t \to 2t$, plugging~\eqref{eq_standar_cont_disc} into the previous display and arguing as in the discrete setting, we conclude that~$G$ has the continuous finite collision property almost surely.
\end{proof}

\end{document}
